% SPDX-License-Identifier: MIT
\usepackage{fontspec}
\setmainfont{texgyrepagella-regular.otf}[
  BoldFont=texgyrepagella-bold.otf,
  ItalicFont=texgyrepagella-italic.otf,
  BoldItalicFont=texgyrepagella-bolditalic.otf]
\setsansfont{texgyreheros-regular.otf}[
  BoldFont=texgyreheros-bold.otf,
  ItalicFont=texgyreheros-italic.otf,
  BoldItalicFont=texgyreheros-bolditalic.otf]
\setmonofont{lmmono10-regular.otf}[Scale=MatchLowercase]
\usepackage[paperwidth=170mm,paperheight=240mm,inner=20mm,outer=18mm,top=20mm,bottom=22mm,headheight=15pt]{geometry}
\usepackage{graphicx,booktabs,tabularx,longtable,array}
\usepackage{amsmath,amssymb}
\usepackage{enumitem}
\usepackage{xcolor}
\definecolor{BookTeal}{HTML}{176C68}
\definecolor{BookInk}{HTML}{23363E}
\definecolor{BookRule}{HTML}{C8D6D6}
\definecolor{BookPale}{HTML}{F1F6F5}
\usepackage{tikz}
\usetikzlibrary{arrows.meta,positioning,shapes.geometric,shapes.misc,fit,backgrounds,calc,matrix}
\usepackage{rotating}
\usepackage{pdfpages}
\usepackage{fvextra}
\DefineVerbatimEnvironment{Prompt}{Verbatim}{fontsize=\small,breaklines=true,breakanywhere=true,frame=single,rulecolor=\color{BookRule},framesep=2mm}
\usepackage{caption}
\captionsetup{font=small,labelfont=bf,labelsep=period}
\usepackage{fancyhdr}
\pagestyle{fancy}
\fancyhf{}
\fancyhead[LE,RO]{\small\thepage}
\fancyhead[LO,RE]{\small\nouppercase{\leftmark}}
\renewcommand{\headrulewidth}{0.3pt}
\usepackage{xurl}
\usepackage[unicode,hypertexnames=false,colorlinks=true,linkcolor=BookTeal,urlcolor=BookTeal,citecolor=BookTeal,breaklinks=true]{hyperref}
\usepackage{bookmark}
\setcounter{tocdepth}{1}
\setcounter{secnumdepth}{1}
\let\cleardoublepage\clearpage
\setlist{nosep,leftmargin=1.8em}
\renewcommand{\arraystretch}{1.25}
\setlength{\parskip}{0.3em}
\setlength{\emergencystretch}{2em}
\newcommand{\bookversion}{0.1.0}
\newcommand{\bookdate}{2026-10-29}
\newcommand{\scidraw}{SciDraw AI}
% Drawn marks: the book fonts lack U+2713/U+2715.
\newcommand{\bookcheck}{\tikz[x=1em,y=1em,baseline=-0.1em,line width=0.12em,line cap=round,line join=round]\draw[color=.] (0,0.3)--(0.28,0)--(0.75,0.7);}
\newcommand{\bookcross}{\tikz[x=1em,y=1em,baseline=-0.1em,line width=0.12em,line cap=round]{\draw[color=.] (0,0)--(0.6,0.65);\draw[color=.] (0,0.65)--(0.6,0);}}
% Circled numbers and the inhibition sign are missing from the book fonts; draw them.
\usepackage{newunicodechar}
\newcommand{\bookcirc}[1]{\tikz[baseline=(n.base)]\node[circle,draw,inner sep=0.5pt,minimum size=1.05em,font=\sffamily\scriptsize] (n) {#1};}
\newcommand{\bookcircfill}[1]{\tikz[baseline=(n.base)]\node[circle,fill,inner sep=0.5pt,minimum size=1.05em,font=\sffamily\scriptsize\bfseries,text=white] (n) {#1};}
\newunicodechar{①}{\bookcirc{1}}
\newunicodechar{②}{\bookcirc{2}}
\newunicodechar{③}{\bookcirc{3}}
\newunicodechar{④}{\bookcirc{4}}
\newunicodechar{⑤}{\bookcirc{5}}
\newunicodechar{⑥}{\bookcirc{6}}
\newunicodechar{⑦}{\bookcirc{7}}
\newunicodechar{⑧}{\bookcirc{8}}
\newunicodechar{⑨}{\bookcirc{9}}
\newunicodechar{⑩}{\bookcirc{10}}
\newunicodechar{⑪}{\bookcirc{11}}
\newunicodechar{⑫}{\bookcirc{12}}
\newunicodechar{⑬}{\bookcirc{13}}
\newunicodechar{❶}{\bookcircfill{1}}
\newunicodechar{❷}{\bookcircfill{2}}
\newunicodechar{❸}{\bookcircfill{3}}
\newunicodechar{❹}{\bookcircfill{4}}
\newunicodechar{❺}{\bookcircfill{5}}
\newunicodechar{❻}{\bookcircfill{6}}
\newunicodechar{❼}{\bookcircfill{7}}
\newunicodechar{⊣}{\ensuremath{\dashv}}
\providecommand{\sdbase}{https://sci-draw.com}
\newcommand{\sdtool}[1]{\href{\sdbase/#1}{\texttt{/#1}}}
\newcommand{\bookrepo}{https://github.com/TopLocalAI/ai-scientific-figures-book}
\newenvironment{BookNote}{\begin{quote}\small\color{BookInk}}{\end{quote}}
\newcolumntype{Y}{>{\raggedright\arraybackslash}X}
\newcommand{\workflowfigure}[2]{%
  \begin{figure}[htbp]\centering
  \input{figures/workflow-#1.tex}
  \caption{#2}\label{fig:workflow}
  \end{figure}}
\newcommand{\pairedfigure}[2]{%
  \begin{figure}[htbp]\centering
  \includegraphics[width=\linewidth]{figures/paired-comparison-#1.pdf}
  \caption{#2}\label{fig:paired}
  \end{figure}}
